\chapter{Sistemas Gênicos}
\label{cap:06}

A célula é uma máquina de informação. Seu genoma codifica não apenas as proteínas que
executam funções estruturais e catalíticas, mas também o programa regulatório que decide
quais proteínas são produzidas, em qual quantidade, em qual momento e em qual tipo celular.
Compreender os sistemas gênicos --- a organização dos genes, os mecanismos de regulação da
expressão e as ferramentas computacionais para acessar e analisar essa informação --- é
condição necessária para a Biologia de Sistemas quantitativa.

Este capítulo cobre o fluxo de informação genética desde a estrutura do DNA até as
proteínas, os diferentes tipos de RNA e suas funções, os mecanismos de regulação da
expressão gênica e sua modelagem matemática, e as principais tecnologias de medição de
expressão --- com ênfase no RNA-seq. Encerramos com as ferramentas computacionais para
localização de genes, acesso a bancos de dados genômicos e análise de enriquecimento
funcional.

\section{O Dogma Central e a Organização dos Genes}

\subsection{O Fluxo de Informação Genética}

\begin{definicaoenv}[Dogma Central da Biologia Molecular]
O dogma central, formulado por Francis Crick em 1958, descreve o fluxo canônico de
informação genética: o DNA é replicado para produzir cópias fiéis de si mesmo; é
transcrito em RNA mensageiro (mRNA); e o mRNA é traduzido em proteína pelo ribossomo.
Em termos esquemáticos: \textbf{DNA} $\xrightarrow{\text{transcrição}}$ \textbf{RNA}
$\xrightarrow{\text{tradução}}$ \textbf{Proteína}.
\end{definicaoenv}

A descoberta da transcriptase reversa por Baltimore e Temin (1970) --- premiada com o Nobel
de 1975 --- revelou que retrovírus como o HIV invertem esse fluxo: convertem seu genoma de
RNA em DNA antes de integrar-se ao genoma do hospedeiro. Esse não é o único complemento ao
dogma: sabe-se hoje que o RNA pode regular outros RNAs (via miRNA e siRNA), que certas
proteínas prions propagam informação estrutural sem envolvimento do DNA, e que a epigenética
transmite informação hereditária por modificações que não alteram a sequência nucleotídica.

Nesse contexto, a palavra ``gene'' ganhou múltiplas definições ao longo do século XX. A
definição operacional mais útil para a Biologia de Sistemas é:

\begin{definicaoenv}[Gene]
Unidade fundamental da hereditariedade: sequência de DNA que codifica uma molécula
funcional --- RNA ou proteína --- e inclui as regiões regulatórias necessárias para
controlar sua expressão (promotor, enhancers, silenciadores).
\end{definicaoenv}

\subsection{Estrutura dos Genes: Elementos Funcionais}

A anatomia de um gene eucariótico típico compreende múltiplos elementos funcionais dispostos
linearmente no DNA. O \emph{promotor} é a região onde a RNA polimerase II e os fatores de
transcrição gerais se ligam para iniciar a transcrição. Em eucariotos, os elementos
promotores mais conservados incluem a TATA box ($\sim$-25 em relação ao sítio de início da
transcrição, consenso TATAAA), a CAAT box ($\sim$-75) e as GC boxes. Reguladores
adicionais, chamados \emph{enhancers} e \emph{silenciadores}, podem atuar a distâncias de
dezenas ou centenas de kilobases do promotor, modulando a atividade transcricional por
looping cromossômico.

A região codificante (ORF, \textit{Open Reading Frame}) é flanqueada pelas regiões não
traduzidas 5' (5' UTR) e 3' (3' UTR). Nos eucariotos, a ORF é interrompida por
\emph{íntrons} --- sequências não codificantes removidas durante o processamento do
pré-mRNA. As sequências codificantes retidas são chamadas \emph{éxons}. O
\emph{splicing alternativo} --- inclusão ou exclusão seletiva de éxons --- permite que um
único gene produza múltiplas isoformas proteicas. O gene da tropomiosina, por exemplo,
produz mais de 60 variantes por splicing alternativo em diferentes tecidos.

Em contraste, os genes procariotos são mais simples: sem íntrons, frequentemente
organizados em \emph{operons} (grupos de genes com função relacionada, transcritos como
um único mRNA policistrónico sob controle de um promotor compartilhado).

\begin{definicaoenv}[Open Reading Frame (ORF)]
Sequência contínua de códons entre um códon de início (ATG, que codifica metionina) e um
códon de parada (TAA, TAG ou TGA), sem códons de parada internos. Uma ORF biologicamente
significativa tipicamente tem comprimento mínimo de 300 nt (100 aminoácidos).
\end{definicaoenv}

\begin{lstlisting}[language=Python, caption=Identificacao de ORFs em sequencia de DNA]
def find_orfs(sequence, min_length=100):
    stop_codons = {"TAA", "TAG", "TGA"}
    orfs = []
    for frame in range(3):
        i = frame
        while i < len(sequence) - 2:
            codon = sequence[i:i+3]
            if codon == "ATG":
                for j in range(i+3, len(sequence)-2, 3):
                    stop = sequence[j:j+3]
                    if stop in stop_codons:
                        if (j - i + 3) >= min_length * 3:
                            orfs.append({'start': i, 'end': j+3,
                                         'frame': frame})
                        break
            i += 3
    return orfs

# Exemplo
seq = "ATGGCTAAATGGCCATAAATG"
print(find_orfs(seq, min_length=1))
\end{lstlisting}


\section{Regulação da Expressão Gênica}

A expressão gênica é regulada em múltiplos níveis: transcricional, pós-transcricional,
traducional e pós-traducional. Cada nível oferece um ponto de controle adicional, permitindo
respostas rápidas, precisas e específicas de contexto. A regulação transcricional --- o
controle da síntese de RNA --- é o nível mais amplamente estudado e o mais amenable à
modelagem matemática quantitativa.

\subsection{Regulação Transcricional em Procariotos: O Operon Lac}

O sistema do operon lac de \textit{Escherichia coli} é o paradigma histórico da regulação
gênica. Proposto por Jacob e Monod em 1961 (Nobel de 1965), o operon lac consiste nos
genes \textit{lacZ} (beta-galactosidase), \textit{lacY} (permeasse) e \textit{lacA}
(transacetilase), todos transcritos sob controle do mesmo promotor. A regulação é dupla:
negativa, pelo repressor LacI, e positiva, pelo ativador CAP.

Na ausência de lactose, o repressor LacI --- uma proteína homodimérica codificada pelo
gene \textit{lacI} adjacente --- liga-se ao operador, uma sequência de DNA sobreposta ao
promotor, bloqueando fisicamente o avanço da RNA polimerase. Quando lactose está presente,
ela é convertida em alolactose, que funciona como indutor: liga-se ao repressor, alterando
sua conformação e reduzindo sua afinidade pelo operador. O repressor se dissocia e a
transcrição ocorre. A regulação positiva pelo CAP (\textit{catabolite activator protein})
garante que o operon lac só seja fortemente ativado na ausência de glicose (a fonte de
carbono preferida), via aumento dos níveis intracelulares de AMPc.

\subsection{Regulação Transcricional em Eucariotos}

Em eucariotos, a regulação transcricional é dramaticamente mais complexa. A cromatina ---
o complexo de DNA e proteínas histônicas --- é o substrato primário da regulação. Genes em
eucromatina (cromatina aberta, acetilada) são transcricionalmente acessíveis; genes em
heterocromatina (cromatina compactada, metilada) são silenciados. As modificações
\emph{epigenéticas} das histonas --- acetilação (H3K27ac: ativação; H3K9me3: repressão) e
metilação do DNA (CpGs metiladas: silenciamento) --- constituem um código adicional de
informação hereditária que não altera a sequência nucleotídica, mas modula
profundamente a expressão gênica.

Os \emph{fatores de transcrição} (TFs) reconhecem sequências regulatórias específicas e
recrutam co-ativadores ou co-repressores para o locus alvo. Ativadores recrutam complexos
remodeladores de cromatina (SWI/SNF) e o Mediador, que conecta TFs ao aparato
transcricional basal. O consórcio ENCODE mapeou sistematicamente elementos regulatórios
no genoma humano por ChIP-seq (imunoprecipitação de cromatina seguida de sequenciamento),
revelando que $\sim$80\% do genoma tem atividade bioquímica mensurável.

\subsection{Modelagem Matemática: A Função de Hill}

A equação de Hill é o modelo matemático padrão para regulação gênica cooperativa. Para
ativação por um fator de transcrição TF:

\begin{equation}
\frac{d[\text{mRNA}]}{dt} = \frac{V_{\max} [\text{TF}]^n}{K_d^n + [\text{TF}]^n} - \delta [\text{mRNA}]
\end{equation}

onde $[\text{TF}]$ é a concentração do fator de transcrição, $n$ é o coeficiente de Hill
(cooperatividade), $K_d$ é a constante de dissociação (concentração de TF que produz
metade da expressão máxima) e $\delta$ é a taxa de degradação do mRNA.

O coeficiente de Hill $n$ determina a ``abrupticidade'' da resposta: para $n = 1$
(Michaelis-Menten), a transição é gradual; para $n \gg 1$, a resposta torna-se quase
binária (\textit{switch-like}), com expressão essencialmente nula abaixo de $K_d$ e máxima
acima. Biologicamente, $n > 1$ indica cooperatividade: múltiplas moléculas de TF ligam-se
cooperativamente ao promotor, de modo que a ocupação do sítio por um TF facilita a ligação
dos demais. O fator de transcrição $p53$, por exemplo, forma um tetrâmero funcional, o que
confere $n \approx 4$ e uma resposta altamente switch-like ao dano de DNA.

Para \emph{repressão} por um TF, a função de Hill é invertida:
\begin{equation}
\frac{d[\text{mRNA}]}{dt} = \frac{V_{\max} K_d^n}{K_d^n + [\text{TF}]^n} - \delta [\text{mRNA}]
\end{equation}

No estado estacionário ($d[\text{mRNA}]/dt = 0$), a concentração de mRNA é:
$[\text{mRNA}]^* = V_{\max} [\text{TF}]^n / [\delta(K_d^n + [\text{TF}]^n)]$.

\begin{lstlisting}[language=Python, caption=Funcao de Hill e estado estacionario]
import numpy as np
import matplotlib.pyplot as plt

def hill_activation(TF, Vmax, Kd, n, delta):
    return (Vmax * TF**n / (Kd**n + TF**n)) / delta

def hill_repression(TF, Vmax, Kd, n, delta):
    return (Vmax * Kd**n / (Kd**n + TF**n)) / delta

TF = np.linspace(0, 50, 300)
Vmax, Kd, delta = 100, 10, 1.0

fig, axes = plt.subplots(1, 2, figsize=(12, 5))
for n in [1, 2, 4, 8]:
    axes[0].plot(TF, hill_activation(TF, Vmax, Kd, n, delta), label=f'n={n}')
    axes[1].plot(TF, hill_repression(TF, Vmax, Kd, n, delta), label=f'n={n}')

for ax, title in zip(axes, ['Ativacao', 'Repressao']):
    ax.axvline(Kd, color='gray', linestyle='--', label=f'Kd={Kd}')
    ax.set_xlabel('[TF]'); ax.set_ylabel('mRNA (estado estacionario)')
    ax.set_title(f'Hill - {title}'); ax.legend(); ax.grid(True)
plt.tight_layout(); plt.show()
\end{lstlisting}


\section{Tipos de RNA e suas Funções}

O transcriptoma é muito mais rico do que o conjunto de mRNAs que codificam proteínas.
Estima-se que $\sim$75\% do genoma humano é transcrito em algum tipo de RNA, mas menos de
2\% codifica proteínas. A vasta maioria dos transcritos são \emph{RNAs não-codificantes}
(ncRNAs) com funções regulatórias, estruturais ou catalíticas.

\subsection{RNA Mensageiro}

O mRNA é o intermediário entre o gene e a proteína. Nos eucariotos, o transcrito primário
(pré-mRNA) sofre um conjunto de modificações co-transcricionais antes de ser exportado para
o citoplasma: adição do \emph{cap} de 7-metilguanosina (m\textsuperscript{7}G) na
extremidade 5', que protege o mRNA da degradação e é reconhecido pela maquinaria traducional;
\emph{splicing} dos íntrons pelo spliceossomo; e poliadenilação na extremidade 3'
(cauda poli-A de $\sim$200 adeninas), que aumenta a estabilidade e facilita a exportação nuclear.

A 5' UTR e a 3' UTR não são regiões inertes: a 5' UTR contém elementos de estrutura
secundária e sequências que regulam a eficiência de iniciação da tradução; a 3' UTR
contém sítios de ligação para miRNAs e proteínas de ligação ao RNA que controlam a
estabilidade e a localização do mRNA. A meia-vida dos mRNAs varia enormemente --- de
minutos (proto-oncogenes como c-Myc e c-Fos) a horas ou dias --- e é um parâmetro
cinético crítico para modelos matemáticos de expressão gênica.

\subsection{RNA de Transferência e RNA Ribossomal}

O tRNA é o adaptador molecular da tradução: carrega o aminoácido adequado e reconhece o
códon correspondente no mRNA via seu anticódon. A estrutura tridimensional em forma de
L do tRNA --- resultado do dobramento da estrutura secundária em ``folha de trevo'' ---
posiciona o aminoácido na extremidade 3' (CCA-OH) enquanto o anticódon interage com o
códon no sítio A do ribossomo. A regra do \emph{wobble} (Crick, 1966) explica por que 61
códons senso são reconhecidos por apenas $\sim$45 tRNAs diferentes: a terceira posição do
anticódon tem flexibilidade de pareamento, podendo reconhecer múltiplos nucletídeos na
terceira posição do códon.

O rRNA é o componente mais abundante e funcional do ribossomo. A subunidade pequena do
ribossomo bacteriano (30S, contendo o 16S rRNA) decoda o mRNA; a subunidade grande (50S,
com 23S e 5S rRNA) catalisa a formação da ligação peptídica. O 16S rRNA é particularmente
importante em bioinformática: sua sequência é altamente conservada mas contém regiões
hipervariáveis que permitem a identificação filogenética de espécies bacterianas ---
base do sequenciamento metagenômico 16S para caracterização de comunidades microbianas.

\subsection{RNAs Regulatórios: miRNA e siRNA}

Os \emph{microRNAs} (miRNAs) são RNAs não-codificantes de $\sim$22 nucleotídeos que
regulam negativamente a expressão de genes alvo por complementaridade parcial com a 3'
UTR dos mRNAs. A biogênese segue uma via multistep: o gene de miRNA é transcrito em
pri-miRNA (centenas de nucleotídeos, com estrutura de hairpin); a endonuclease Drosha
(no núcleo) processa o pri-miRNA em pré-miRNA ($\sim$70 nt); a Exportina-5 transporta o
pré-miRNA para o citoplasma; a Dicer gera o duplex de miRNA maduro ($\sim$22 nt); e uma
fita do duplex é incorporada ao complexo RISC (\textit{RNA-Induced Silencing Complex}).
O RISC usa o miRNA como guia para encontrar mRNAs complementares: pareamento perfeito
induz clivagem do mRNA; pareamento imperfeito (mais comum em animais) induz repressão
traducional e deadenilação.

A importância biológica dos miRNAs é imensa: um único miRNA pode regular centenas de
mRNAs, e o genoma humano codifica $>2500$ miRNAs (miRBase). Alterações na expressão de
miRNAs estão associadas a praticamente todos os tipos de câncer, e miRNAs são
investigados como biomarcadores diagnósticos no plasma sanguíneo.

Os \emph{siRNAs} (\textit{small interfering RNAs}) têm estrutura similar aos miRNAs mas
atuam por pareamento perfeito, induzindo clivagem do mRNA alvo com alta especificidade.
A interferência por RNA (RNAi), descoberta por Fire e Mello em 1998 (Nobel 2006), é hoje
uma ferramenta essencial para knockdown gênico experimental e está em desenvolvimento como
plataforma terapêutica.

Os \emph{lncRNAs} (\textit{long non-coding RNAs}) têm mais de 200 nt e funções diversas:
regulação da estrutura da cromatina em \textit{cis} (HOTAIR, XIST para inativação do X) ou
em \textit{trans}, atuando como esponjas de miRNAs (ceRNAs) ou como scaffolds para
complexos proteicos.


\section{Medição da Expressão Gênica}

\subsection{Da Hibridização ao Sequenciamento}

A medição sistemática da expressão gênica evoluiu rapidamente nas últimas décadas. O
\emph{Northern blot} (1977) permitiu detectar transcritos individuais por hibridização com
sondas marcadas, mas processa apenas um gene por vez. O \emph{RT-PCR quantitativo} (qPCR)
é o padrão-ouro para validação de expressão de genes específicos: o mRNA é convertido em
cDNA por transcriptase reversa, e o cDNA é amplificado em presença de fluoróforos que
permitem quantificação em tempo real.

Os \emph{microarrays de DNA} (1995, Brown e Botstein) revolucionaram a biologia ao permitir
a medição simultânea da expressão de milhares de genes. Sondas de DNA (cDNA ou
oligonucleotídeos) são imobilizadas em chips de vidro ou silício; o cDNA marcado com
fluoróforos derivado das amostras experimental e controle é hibridizado ao chip; e a
intensidade de fluorescência de cada spot reflete a abundância relativa de cada transcrito.
Apesar das limitações --- faixa dinâmica restrita, necessidade de sondas pré-definidas,
hibridização cruzada --- os microarrays geraram dados que ainda são amplamente estudados
nos repositórios públicos GEO (Gene Expression Omnibus) e ArrayExpress.

\subsection{RNA-seq}

O \emph{RNA-seq} (sequenciamento de RNA por sequenciadores de alto rendimento) tornou-se
rapidamente o método de referência para análise do transcriptoma, suplantando os
microarrays. No protocolo padrão, o RNA total (ou o RNA poli-A-selecionado) é fragmentado,
convertido em cDNA de fita dupla e preparado numa biblioteca de fragmentos de tamanho
controlado. Após sequenciamento (Illumina: 50--300 nt por read), os reads são alinhados ao
genoma de referência e contabilizados por gene, produzindo uma matriz de contagens
(genes × amostras).

As vantagens do RNA-seq sobre os microarrays são substanciais: (i) faixa dinâmica de cinco
ordens de magnitude, versus três nos microarrays; (ii) ausência de sondas pré-definidas,
permitindo descoberta de novos transcritos e genes; (iii) resolução de variantes de
splicing; (iv) quantificação absoluta em reads por gene; e (v) compatibilidade com
amostras de qualidade variável (tumor formalizado, sangue seco). As desvantagens incluem o
custo computacional de análise e a necessidade de um genoma de referência de qualidade para
o alinhamento.

O pipeline padrão de RNA-seq segue as etapas: (1) controle de qualidade dos reads (FastQC);
(2) trimagem de adaptadores e bases de baixa qualidade (Trimmomatic, fastp); (3)
alinhamento ao genoma de referência (STAR, HISAT2); (4) contagem de reads por gene
(featureCounts, HTSeq); (5) normalização e análise de expressão diferencial (DESeq2,
edgeR); (6) análise de enriquecimento funcional (clusterProfiler, GSEA).

\subsection{Normalização e Expressão Diferencial}

A normalização de dados de RNA-seq visa remover variações técnicas (número de reads por
amostra, comprimento dos genes) para tornar amostras comparáveis. As métricas mais comuns
são:

\begin{itemize}
\item \textbf{RPKM/FPKM} (\textit{Reads/Fragments Per Kilobase per Million}): normaliza
  por comprimento do gene e pelo número total de reads. Inadequado para comparações entre
  amostras porque o total de RPKM não é constante.
\item \textbf{TPM} (\textit{Transcripts Per Million}): normaliza primeiro por comprimento,
  depois pelo total. Comparável entre amostras do mesmo tecido.
\item \textbf{DESeq2 VST} e \textbf{TMM} (edgeR): normalização baseada em modelos
  estatísticos negativos binomiais, preferida para análise de expressão diferencial formal.
\end{itemize}

A análise de expressão diferencial identifica genes cuja expressão difere significativamente
entre condições (tratado vs.\ controle, tumor vs.\ normal). O DESeq2 modela as contagens
com distribuição binomial negativa, estima dispersões de forma shrunk (emprestando força
estatística entre genes similares) e aplica o teste de Wald para calcular $p$-valores. A
correção de Benjamini-Hochberg (FDR) controla a taxa de falsos positivos no contexto de
múltiplos testes simultâneos.

\begin{lstlisting}[language=Python, caption=Pipeline de RNA-seq e volcano plot]
import numpy as np
import pandas as pd
import matplotlib.pyplot as plt
from scipy import stats

# Simular matriz de contagens (30 genes, 6 amostras: 3 controle + 3 tratamento)
np.random.seed(42)
n_genes = 500
ctrl = np.random.negative_binomial(50, 0.5, (n_genes, 3))
trt  = np.random.negative_binomial(50, 0.5, (n_genes, 3))
# 50 genes diferencialmente expressos
trt[:50] = np.random.negative_binomial(200, 0.5, (50, 3))
counts = np.hstack([ctrl, trt])

# Normalizacao TPM simplificada
gene_lengths = np.random.randint(500, 5000, n_genes)
def tpm(counts, lengths):
    rpk = counts / (lengths[:, None] / 1000)
    sf = rpk.sum(axis=0) / 1e6
    return rpk / sf

tpm_counts = tpm(counts, gene_lengths)

# Teste t simples por gene
log_fc = np.log2(tpm_counts[:, 3:].mean(1) / (tpm_counts[:, :3].mean(1) + 1e-6))
p_vals = np.array([stats.ttest_ind(tpm_counts[i, :3], tpm_counts[i, 3:])[1]
                   for i in range(n_genes)])

# Correcao de Benjamini-Hochberg
from statsmodels.stats.multitest import multipletests
_, p_adj, _, _ = multipletests(p_vals, method='fdr_bh')

# Volcano plot
sig = (p_adj < 0.05) & (np.abs(log_fc) > 1)
plt.figure(figsize=(8, 6))
plt.scatter(log_fc, -np.log10(p_adj), alpha=0.4, s=10, color='gray')
plt.scatter(log_fc[sig], -np.log10(p_adj[sig]), color='red', s=15,
            label=f'Significativos: {sig.sum()}')
plt.axhline(-np.log10(0.05), color='r', linestyle='--')
plt.axvline(-1, color='b', linestyle='--'); plt.axvline(1, color='b', linestyle='--')
plt.xlabel('log2(Fold Change)'); plt.ylabel('-log10(FDR)')
plt.title('Volcano Plot --- RNA-seq'); plt.legend(); plt.grid(True, alpha=0.3)
plt.tight_layout(); plt.show()
\end{lstlisting}


\section{Localização e Anotação de Genes}

\subsection{Predição de Genes}

A identificação de genes a partir da sequência genômica crua é um dos problemas fundacionais
da bioinformática. As abordagens se dividem em dois grandes paradigmas. Na predição
\emph{ab initio}, algoritmos identificam genes com base em características intrínsecas da
sequência: presença de ORFs, composição de dinucleotídeos, modelos de sítios de splice
(doador GT e aceitador AG) e sinais de poliadenilação. Os modelos de Markov Ocultos (HMMs)
são a arquitetura padrão para esse tipo de predição --- Augustus e GenScan (eucariotos),
Prodigal e GeneMark (procariotos) são as ferramentas mais usadas.

Na predição baseada em \emph{evidências}, dados experimentais orientam a anotação:
alinhamentos de ESTs (\textit{Expressed Sequence Tags}) e transcritos completos de cDNA,
proteínas homólogas de espécies relacionadas por BLAST, e reads de RNA-seq mapeados ao
genoma. Pipelines integrativos como MAKER e Braker combinam as duas abordagens para
produzir anotações de alta qualidade.

A predição de genes em eucariotos é significativamente mais difícil que em procariotos,
pelos seguintes motivos: (i) os genes eucarióticos são esparsos --- $\sim$25\,000 genes em
$\sim$3 bilhões de pares de bases humanos; (ii) os íntrons são longos (mediana $\sim$3 kb)
e numerosos (média de 8 éxons por gene humano); (iii) o splicing alternativo multiplica as
isoformas; e (iv) o DNA repetitivo ($\sim$45\% do genoma humano) interfere com os
alinhamentos.

\subsection{Bancos de Dados Genômicos}

O GenBank (NCBI) é o repositório primário de sequências de nucleotídeos, com mais de 350
milhões de sequências e 10 trilhões de pares de bases (2024). Cada registro genômico inclui
a sequência, metadados do organismo e uma tabela de \emph{features} anotando elementos
funcionais: genes, CDSs, mRNAs, rRNAs, repeat regions. O formato GenBank é texto plano
estruturado com campos padronizados (LOCUS, DEFINITION, ACCESSION, FEATURES, ORIGIN).

O BioPython (Cock et al., 2009) fornece uma interface Python completa para acessar e
manipular registros do NCBI via Entrez:

\begin{lstlisting}[language=Python, caption=Acesso ao GenBank via BioPython]
from Bio import Entrez, SeqIO

Entrez.email = "seu_email@example.com"

def fetch_gene(query, db="nucleotide", max_results=5):
    handle = Entrez.esearch(db=db, term=query, retmax=max_results)
    ids = Entrez.read(handle)["IdList"]
    records = []
    for gid in ids:
        h = Entrez.efetch(db=db, id=gid, rettype="gb", retmode="text")
        rec = SeqIO.read(h, "genbank")
        records.append(rec)
    return records

# Buscar HBB (hemoglobina beta humana)
recs = fetch_gene("Homo sapiens HBB mRNA", max_results=1)
rec = recs[0]
print(f"ID: {rec.id} | {len(rec)} bp | {rec.annotations['organism']}")

# Extrair CDSs
for feat in rec.features:
    if feat.type == "CDS":
        prod = feat.qualifiers.get("product", ["?"])[0]
        prot = feat.qualifiers.get("translation", [""])[0]
        print(f"CDS: {prod} | {len(prot)} aa")
\end{lstlisting}

\subsection{Ontologia Gênica e Análise de Enriquecimento}

A Gene Ontology (GO) é um vocabulário controlado que descreve funções de genes em três
domínios hierárquicos: \textbf{Função Molecular} (MF: atividades bioquímicas da proteína,
ex: ``oxygen transporter activity''), \textbf{Processo Biológico} (BP: processos de nível
celular ou organísmico, ex: ``oxygen transport'') e \textbf{Componente Celular} (CC:
localização, ex: ``hemoglobin complex''). Os termos GO formam um DAG (\textit{Directed
Acyclic Graph}) em que termos mais específicos são filhos de termos mais gerais.

A análise de \emph{enriquecimento de GO} responde à pergunta: dado um conjunto de genes
(por exemplo, genes diferencialmente expressos num experimento de RNA-seq), quais termos
GO estão sobre-representados em comparação com o genoma de referência? O teste mais comum
é o hipergeométrico (ou equivalentemente, o teste exato de Fisher), com correção para
múltiplos testes por FDR. Ferramentas como clusterProfiler (R), gProfiler e Enrichr
automatizam essa análise e produzem visualizações como dotplots e redes de termos GO.

Além da GO, as análises de enriquecimento frequentemente incluem vias do KEGG
(\textit{Kyoto Encyclopedia of Genes and Genomes}) e do Reactome, que organizam genes em
vias metabólicas e de sinalização biologicamente interpretáveis.


\section{Exercícios}

\begin{exercicio}
Uma região codificante tem 1200 nucleotídeos. (a) Quantos códons ela contém? (b) Quantos
aminoácidos codifica a proteína resultante (descontando o códon de parada)? (c) Se a taxa
de tradução em eucariotos é de $\sim$5 aminoácidos por segundo, quanto tempo leva a
síntese completa de uma molécula dessa proteína?
\end{exercicio}

\begin{exercicio}
Calcule o conteúdo GC da sequência \texttt{5'-ATGCGATCGTAGCTAGCTAA-3'} e discuta por que
o conteúdo GC é biologicamente relevante: como ele afeta a temperatura de desnaturação
do DNA ($T_m$) e por que regiões promotoras de genes ativamente transcritos tendem a ter
conteúdo GC baixo?
\end{exercicio}

\begin{exercicio}
Para uma função de Hill de ativação com $n = 4$, $K_d = 10$ nM e $V_{\max} = 100$
nM/min: (a) calcule a taxa de transcrição quando $[\text{TF}] = 10$ nM; (b) em que
concentração de TF se atinge 90\% de $V_{\max}$? (c) compare com o caso $n = 1$ ---
qual requer maior variação de $[\text{TF}]$ para passar de 10\% a 90\% de $V_{\max}$?
Interprete biologicamente o papel do coeficiente de Hill na robustez do switch.
\end{exercicio}

\begin{exercicio}[RNA-seq]
Implemente em Python uma função que: (a) normalize uma matriz de contagens brutas para
TPM, dado um vetor de comprimentos de genes; (b) calcule log2(Fold Change) e valor-$p$
(teste $t$ de Welch) entre dois grupos de amostras; (c) aplique correção de Benjamini-Hochberg
para FDR; (d) gere um volcano plot colorindo em vermelho os genes com FDR $<$ 0.05 e
$|\text{log2FC}| > 1$.
\end{exercicio}

\begin{exercicio}[BioPython]
Usando BioPython: (a) baixe o registro GenBank do gene TP53 humano (busca: ``Homo sapiens
TP53 mRNA RefSeq''); (b) extraia e imprima todas as features do tipo CDS e gene;
(c) calcule o conteúdo GC da sequência codificante; (d) traduza a CDS e verifique quantos
aminoácidos tem a proteína p53.
\end{exercicio}

\begin{exercicio}
Considere o modelo de regulação gênica com síntese por função de Hill e degradação linear:
\begin{align*}
  \frac{dm}{dt} &= \frac{V_m [P]^n}{K_m^n + [P]^n} - \delta_m m \\
  \frac{d[P]}{dt} &= k_p m - \delta_p [P]
\end{align*}
onde $m$ é o mRNA e $[P]$ é a proteína. Com $V_m = 10$, $K_m = 5$, $n = 2$,
$\delta_m = 0.5$, $k_p = 2$, $\delta_p = 0.3$: (a) encontre o estado estacionário
numericamente; (b) simule as dinâmicas por 50 unidades de tempo a partir de $(m_0, P_0) = (0, 0)$;
(c) analise a estabilidade do equilíbrio via jacobiano.
\end{exercicio}

\begin{exercicio}[Projeto Integrador]
Escolha um conjunto de dados de RNA-seq disponível no GEO (ex: GSE74246, carcinoma
hepático). Construa um pipeline completo: (a) baixe os dados de contagem pré-processados;
(b) normalize e filtre genes com baixa contagem; (c) realize análise de expressão
diferencial entre dois grupos biológicos; (d) gere volcano plot e heatmap dos 50 genes
mais significativos; (e) realize análise de enriquecimento GO e KEGG com os genes
diferencialmente expressos; (f) interprete biologicamente os resultados: quais vias estão
alteradas? Há consistência com a literatura?
\end{exercicio}

\section{Notas Bibliográficas}

O dogma central da biologia molecular foi enunciado por Crick (1958) \cite{crick1958} e
formalizado em detalhes em Crick (1970) \cite{crick1970}. A descoberta da transcriptase
reversa por Baltimore \cite{baltimore1970} e Temin e Mizutani \cite{temin1970} (Nobel 1975)
representa a primeira exceção demonstrada ao fluxo canônico. A interferência por RNA foi
descoberta por Fire e Mello \cite{fire1998} (Nobel 2006). O consórcio ENCODE mapeou
sistematicamente os elementos regulatórios do genoma humano \cite{encode2012}.

Para biologia molecular clássica, Watson et al.\ \cite{watson2013} (\textit{Molecular
Biology of the Gene}) e Alberts et al.\ \cite{alberts2022} (\textit{Molecular Biology of
the Cell}) são as referências textuais definitivas. O BioPython é descrito em Cock et al.\
\cite{cock2009}. DESeq2 é descrito em Love et al.\ \cite{love2014} e é a ferramenta de
análise de expressão diferencial de RNA-seq mais citada. Para microRNAs, Bartel (2018)
\cite{bartel2018} fornece uma revisão abrangente.

Para a função de Hill e modelagem matemática de regulação gênica, o livro de Alon (2019)
\cite{alon2019} é a referência mais acessível e matematicamente rigorosa. O RNA-seq é
introduzido em Wang et al.\ (2009) \cite{wang2009} e revisado anualmente com novas
ferramentas.

\begin{thebibliography}{99}

\bibitem{crick1958}
Crick F.H.C. (1958).
On protein synthesis.
\textit{Symposia of the Society for Experimental Biology} \textbf{12}: 138--163.

\bibitem{crick1970}
Crick F.H.C. (1970).
Central dogma of molecular biology.
\textit{Nature} \textbf{227}: 561--563.

\bibitem{baltimore1970}
Baltimore D. (1970).
RNA-dependent DNA polymerase in virions of RNA tumour viruses.
\textit{Nature} \textbf{226}: 1209--1211.

\bibitem{temin1970}
Temin H.M., Mizutani S. (1970).
RNA-dependent DNA polymerase in virions of Rous sarcoma virus.
\textit{Nature} \textbf{226}: 1211--1213.

\bibitem{fire1998}
Fire A. et al. (1998).
Potent and specific genetic interference by double-stranded RNA in
\textit{Caenorhabditis elegans}.
\textit{Nature} \textbf{391}: 806--811.

\bibitem{encode2012}
ENCODE Project Consortium (2012).
An integrated encyclopedia of DNA elements in the human genome.
\textit{Nature} \textbf{489}: 57--74.

\bibitem{watson2013}
Watson J.D. et al. (2013).
\textit{Molecular Biology of the Gene}, 7\textordfeminine\ ed.
Pearson.

\bibitem{alberts2022}
Alberts B. et al. (2022).
\textit{Molecular Biology of the Cell}, 7\textordfeminine\ ed.
Garland Science.

\bibitem{cock2009}
Cock P.J.A. et al. (2009).
Biopython: freely available Python tools for computational molecular biology
and bioinformatics.
\textit{Bioinformatics} \textbf{25}: 1422--1423.

\bibitem{love2014}
Love M.I., Huber W., Anders S. (2014).
Moderated estimation of fold change and dispersion for RNA-seq data with DESeq2.
\textit{Genome Biology} \textbf{15}: 550.

\bibitem{bartel2018}
Bartel D.P. (2018).
Metazoan MicroRNAs.
\textit{Cell} \textbf{173}: 20--51.

\bibitem{alon2019}
Alon U. (2019).
\textit{An Introduction to Systems Biology: Design Principles of Biological Circuits},
2\textordfeminine\ ed.
Chapman \& Hall/CRC.

\bibitem{wang2009}
Wang Z., Gerstein M., Snyder M. (2009).
RNA-Seq: a revolutionary tool for transcriptomics.
\textit{Nature Reviews Genetics} \textbf{10}: 57--63.

\end{thebibliography}
